\input{preamble}

% OK, start here.
%
\begin{document}

\title{Théorie des ensembles}


\maketitle

\phantomsection
\label{section-phantom}

\tableofcontents

\section{Introduction}
\label{section-introduction}

\noindent
Nous avons besoin de quelques notions de théorie des ensembles de temps à autre.
Nous utilisons la théorie des ensembles de Zermelo--Fraenkel avec l'axiome du
choix (ZFC), telle qu'elle est décrite dans \cite{Kunen} et \cite{Jech}.

\section{Tout est un ensemble}
\label{section-sets-everything}

\noindent
La plupart des mathématiciens considèrent que la théorie des ensembles fournit
les fondements de base des mathématiques. Mais comment cela fonctionne-t-il
réellement ? Par exemple, comment traduit-on la phrase
« $X$ est un schéma » en théorie des ensembles ? Il suffit de dérouler les
définitions : un schéma est un espace localement annelé tel que tout point
possède un voisinage ouvert qui est un schéma affine. Un espace localement
annelé est un espace annelé en anneaux locaux. Un espace annelé est une paire
$(X, \mathcal{O}_X)$ constituée
d'un espace topologique $X$ et d'un faisceau d'anneaux $\mathcal{O}_X$ sur
celui-ci. Un espace topologique est une paire $(X, \tau)$ constituée d'un
ensemble $X$ et d'un ensemble de sous-ensembles
$\tau \subset \mathcal{P}(X)$ satisfaisant les axiomes d'une topologie. Et
ainsi de suite.

\medskip\noindent
Ainsi, étant donné un ensemble $S$, comment reconnaîtrions-nous s'il s'agit
d'un schéma ? La première chose que nous cherchons à déterminer est si
l'ensemble $S$ est une paire ordonnée. Cela signifie par définition (voir
\cite{Jech}, page 7) que $S$ est de la forme
$(a, b) := \{\{a\}, \{a, b\}\}$ pour certains ensembles $a, b$. Si c'est le
cas, nous examinerions alors si $a$ est une paire ordonnée $(c, d)$. Si oui,
nous vérifierions si $d \subset \mathcal{P}(c)$, puis, dans l'affirmative, si
$d$ constitue la collection d'ensembles d'une topologie sur l'ensemble $c$.
Et ainsi de suite.

\medskip\noindent
Ainsi, même si la rédaction d'une formule complète $\phi_{scheme}(x)$ à une
variable libre $x$ en théorie des ensembles, exprimant la notion « $x$ est un
schéma », demanderait un travail considérable, elle est possible. Il devrait en
être de même pour tout objet mathématique.

\section{Classes}
\label{section-classes}

\noindent
Nous utilisons informellement la notion de {\it classe}.\par\noindent
Étant donnée une formule $\phi(x, p_1, \ldots, p_n)$, nous appelons
$$
C = \{x : \phi(x, p_1, \ldots, p_n)\}
$$
une {\it classe}. Une classe est plus facile à manipuler que la formule qui la
définit, mais ce n'est pas, à proprement parler, un objet mathématique. Par
exemple, si $R$ est un anneau, nous pouvons considérer la classe de tous les
$R$-modules (puisque, après tout, nous pouvons traduire la phrase « $M$ est un
$R$-module » en une formule de théorie des ensembles, laquelle définit alors
une classe). Une {\it classe propre} est une classe qui n'est pas un ensemble.

\medskip\noindent
Nous pouvons ainsi considérer la catégorie des $R$-modules, qui est une
« grande » catégorie---autrement dit, elle possède une classe propre d'objets.
De même, nous pouvons considérer la « grande » catégorie des schémas, la
« grande » catégorie des anneaux, etc.



\section{Ordinaux}
\label{section-ordinals}

\noindent
Un ensemble $T$ est {\it transitif} si $x\in T$ implique $x\subset T$.
Un ensemble $\alpha$ est un {\it ordinal} s'il est transitif et bien ordonné
par $\in$. Dans ce cas, nous définissons
$\alpha + 1 = \alpha \cup \{\alpha\}$, qui est un autre ordinal appelé le
{\it successeur} de $\alpha$. Un ordinal $\alpha$ est appelé un
{\it ordinal successeur} s'il existe un ordinal $\beta$ tel que
$\alpha = \beta + 1$. Le plus petit ordinal est $\emptyset$, également noté
$0$. Si $\alpha$ n'est ni $0$ ni un ordinal successeur, alors $\alpha$ est
appelé un {\it ordinal limite} et nous avons
$$
\alpha
=
\bigcup\nolimits_{\gamma \in \alpha} \gamma.
$$
Le premier ordinal limite est $\omega$ ; c'est aussi le premier ordinal
infini. Le premier ordinal non dénombrable $\omega_1$ est l'ensemble de tous
les ordinaux dénombrables. La collection de tous les ordinaux est une classe
propre. Elle est bien ordonnée par $\in$ au sens suivant : tout ensemble non
vide (et même toute classe non vide) d'ordinaux possède un plus petit élément.
Étant donné un ensemble $A$ d'ordinaux, nous définissons la {\it borne supérieure} de
$A$ par $\sup_{\alpha \in A} \alpha =
\bigcup_{\alpha \in A} \alpha$. C'est le plus petit ordinal supérieur ou égal
à tout $\alpha \in A$. Étant donné un ensemble bien ordonné $(S, <)$, il existe
un unique ordinal $\alpha$ tel que $(S, <) \cong (\alpha, \in)$ ; on l'appelle
le {\it type d'ordre} de l'ensemble bien ordonné.

\section{La hiérarchie des ensembles}
\label{section-sets-hierarchy}

\noindent
Nous définissons par récurrence transfinie $V_0 = \emptyset$,
$V_{\alpha + 1} = P(V_\alpha)$ (ensemble des parties), et, pour un ordinal
limite $\alpha$,
$$
V_\alpha = \bigcup\nolimits_{\beta < \alpha} V_\beta.
$$
Remarquons que chaque $V_\alpha$ est un ensemble transitif.

\begin{lemma}
\label{lemma-axiom-regularity}
Tout ensemble est un élément de $V_\alpha$ pour un certain ordinal $\alpha$.
\end{lemma}

\begin{proof}
Voir \cite[lemme 6.3]{Jech}.
\end{proof}

\noindent
Dans \cite[chapitre III]{Kunen}, il est expliqué que ce lemme équivaut à
l'axiome de fondation. Le {\it rang} d'un ensemble $S$ est le plus petit ordinal
$\alpha$ tel que $S \in V_{\alpha + 1}$. Par {\it univers partiel}, nous
entendrons un ensemble suffisamment grand de la forme $V_\alpha$, lequel sera
clair d'après le contexte.

\section{Cardinalité}
\label{section-cardinals}

\noindent
La {\it cardinalité} d'un ensemble $A$ est le plus petit ordinal $\alpha$ tel
qu'il existe une bijection entre $A$ et $\alpha$. Nous utilisons parfois la
notation $\alpha = |A|$ pour l'indiquer. Nous disons qu'un ordinal $\alpha$ est
un {\it cardinal} si et seulement s'il apparaît comme cardinalité d'un certain
ensemble $A$---autrement dit, si $\alpha = |A|$. Nous utilisons les lettres
grecques $\kappa$, $\lambda$ pour les cardinaux. Le premier cardinal infini est
$\omega$ et, dans ce contexte, il est noté $\aleph_0$. Un ensemble est
{\it dénombrable} si sa cardinalité est $\leq \aleph_0$. Si $\alpha$ est un
ordinal, nous notons $\alpha^+$ le plus petit cardinal $> \alpha$. On peut
l'utiliser pour définir $\aleph_1 = \aleph_0^+$,
$\aleph_2 = \aleph_1^+$, etc., et l'on peut en fait définir $\aleph_\alpha$
pour tout ordinal $\alpha$ par récurrence transfinie. Nous relevons l'égalité
$\aleph_1 = \omega_1$.

\medskip\noindent
L'{\it addition} des cardinaux $\kappa, \lambda$ est notée
$\kappa \oplus \lambda$ ; c'est la cardinalité de
$\kappa \amalg \lambda$. La {\it multiplication} des cardinaux
$\kappa, \lambda$ est notée $\kappa \otimes \lambda$ ; c'est la cardinalité
de $\kappa \times \lambda$. Si $\kappa$ et $\lambda$ sont des cardinaux
infinis, alors
$\kappa \oplus \lambda = \kappa \otimes \lambda = \max(\kappa, \lambda)$.
L'{\it exponentiation} des cardinaux $\kappa, \lambda$ est notée
$\kappa^\lambda$ ; c'est la cardinalité de l'ensemble des applications
(ensemblistes) de $\lambda$ vers $\kappa$. Étant donné un ensemble $K$ de
cardinaux, la {\it borne supérieure} de $K$ est
$\sup_{\kappa \in K} \kappa = \bigcup_{\kappa \in K} \kappa$, qui est
également un cardinal.

\section{Cofinalité}
\label{section-cofinality}

\noindent
Un {\it sous-ensemble cofinal} $S$ d'un ensemble bien ordonné $T$ est un
sous-ensemble $S \subset T$ tel que
$\forall t \in T \exists s\in S (t \leq s)$. Remarquons qu'un sous-ensemble
d'un ensemble bien ordonné est lui-même bien ordonné (pour l'ordre induit).
Étant donné un ordinal $\alpha$, la {\it cofinalité} $\text{cf}(\alpha)$ de
$\alpha$ est le plus petit ordinal $\beta$ qui apparaît comme type d'ordre
d'un sous-ensemble cofinal de $\alpha$. La cofinalité d'un ordinal est toujours
un cardinal. Nous pouvons donc aussi définir la cofinalité de $\alpha$ comme la
plus petite cardinalité d'un sous-ensemble cofinal de $\alpha$.

\begin{lemma}
\label{lemma-map-from-set-lifts}
Supposons que $T = \colim_{\alpha < \beta} T_\alpha$ soit une colimite
d'ensembles indexée par les ordinaux strictement inférieurs à un ordinal donné
$\beta$. Supposons que $\varphi : S \to T$ soit une application d'ensembles.
Alors $\varphi$ se relève en une application à valeurs dans $T_\alpha$ pour un
certain $\alpha < \beta$, pourvu que $\beta$ ne soit pas la limite d'une
famille d'ordinaux indexée par $S$, autrement dit si $\beta$ est un ordinal tel que
$\text{cf}(\beta) > |S|$.
\end{lemma}

\begin{proof}
Pour chaque élément $s \in S$, choisissons un $\alpha_s < \beta$ et un élément
$t_s \in T_{\alpha_s}$ qui s'envoie sur $\varphi(s)$ dans $T$. Par hypothèse,
$\alpha = \sup_{s \in S} \alpha_s$ est strictement inférieur à $\beta$.
L'application $\varphi_\alpha : S \to T_\alpha$ qui associe à $s$ l'image de
$t_s$ dans $T_\alpha$ est donc une solution.
\end{proof}

\noindent
Le résultat suivant est essentiellement l'argument de Grothendieck établissant
l'existence d'ordinaux de cofinalité arbitrairement grande, argument qu'il a
utilisé pour démontrer l'existence d'assez d'injectifs dans certaines
catégories abéliennes ; voir \cite{Tohoku}.

\begin{proposition}
\label{proposition-exist-ordinals-large-cofinality}
Soit $\kappa$ un cardinal. Il existe alors un ordinal dont la cofinalité est
strictement supérieure à $\kappa$.
\end{proposition}

\begin{proof}
Si $\kappa$ est fini, alors $\omega = \text{cf}(\omega)$ convient. Supposons
donc $\kappa$ infini. Considérons le plus petit ordinal $\alpha$ dont la
cardinalité est strictement supérieure à $\kappa$. Nous affirmons que
$\text{cf}(\alpha) > \kappa$. Remarquons que $\alpha$ est un ordinal limite,
car si $\alpha = \beta + 1$, alors $|\alpha| = |\beta|$ (puisque $\alpha$ et
$\beta$ sont infinis), ce qui contredit la minimalité de $\alpha$. (Bien sûr,
$\alpha$ est aussi un cardinal, mais nous n'en avons pas besoin.) Pour obtenir
une contradiction, supposons que $S \subset \alpha$ soit un sous-ensemble
cofinal tel que $|S| \leq \kappa$. Pour $\beta \in S$, c'est-à-dire
$\beta < \alpha$, nous avons $|\beta| \leq \kappa$ par minimalité de $\alpha$.
Comme $\alpha$ est un ordinal limite et $S$ est cofinal dans $\alpha$, nous
obtenons $\alpha = \bigcup_{\beta \in S} \beta$. Par conséquent,
$|\alpha| \leq |S| \otimes \kappa \leq \kappa \otimes \kappa \leq \kappa$,
ce qui contredit notre choix de $\alpha$.
\end{proof}



\section{Principe de réflexion}
\label{section-reflection-principle}

\noindent
Une partie de ce contenu se trouve dans le chapitre de \cite{Kunen} intitulé
« Easy consistency proofs ».

\medskip\noindent
Soit $\phi(x_1, \ldots, x_n)$ une formule de théorie des ensembles. Adoptons la
convention selon laquelle cette notation implique que toutes les variables
libres de $\phi$ figurent parmi $x_1, \ldots, x_n$. Soit $M$ un ensemble. La
formule $\phi^M(x_1, \ldots, x_n)$ est celle que l'on obtient à partir de
$\phi(x_1, \ldots, x_n)$ en remplaçant respectivement tous les $\forall x$ et
$\exists x$ par $\forall x\in M$ et $\exists x\in M$. Ainsi, la formule
$\phi(x_1, x_2) = \exists x (x\in x_1 \wedge x\in x_2)$ devient
$\phi^M(x_1, x_2) = \exists x \in M (x\in x_1 \wedge x\in x_2)$.
La formule $\phi^M$ est appelée la {\it relativisation de $\phi$ à $M$}.

\begin{theorem}
\label{theorem-reflection-principle}
Considérons une famille {\bf finie} de formules de théorie des ensembles,
notées $\phi_1(x_1, \ldots, x_n), \ldots, \phi_m(x_1, \ldots, x_n)$.
Soit $M_0$ un ensemble. Il existe un ensemble $M$ tel
que $M_0 \subset M$ et que l'on ait, $\forall x_1, \ldots, x_n \in M$,
$$
\forall i = 1, \ldots, m, \ \phi_i^{M}(x_1, \ldots, x_n)
\Leftrightarrow
\forall i = 1, \ldots, m, \ \phi_i(x_1, \ldots, x_n).
$$
En fait, nous pouvons prendre $M = V_\alpha$ pour un certain ordinal limite
$\alpha$.
\end{theorem}

\begin{proof}
Voir \cite[théorème 12.14]{Jech} ou \cite[théorème 7.4]{Kunen}.
\end{proof}

\noindent
Nous interprétons ce théorème comme suit : étant donnés
$x_1, \ldots, x_n \in M$, les formules sont vraies lorsque les variables liées
parcourent tous les ensembles si et seulement si elles sont vraies lorsque les
variables liées parcourent les éléments de $V_\alpha$. Ce théorème est un
métathéorème, car il porte directement sur les formules de la théorie des
ensembles. Il affirme en réalité qu'étant donnée la liste finie de formules
$\phi_1, \ldots, \phi_m$, dont toutes les variables libres figurent parmi
$x_1, \ldots, x_n$, la phrase
$$
\begin{matrix}
\forall M_0\ \exists M, \ M_0 \subset M\ \forall x_1, \ldots, x_n \in M \\
\phi_1(x_1, \ldots, x_n) \wedge \ldots \wedge \phi_m(x_1, \ldots, x_n)
\leftrightarrow
\phi_1^M(x_1, \ldots, x_n) \wedge \ldots \wedge \phi_m^M(x_1, \ldots, x_n)
\end{matrix}
$$
est démontrable dans ZFC. Autrement dit, chaque fois que nous écrivons
effectivement une liste finie de formules $\phi_i$, nous obtenons un théorème.

\medskip\noindent
Il est assez difficile d'utiliser ce théorème dans les « mathématiques
ordinaires », car le sens des formules $\phi_i^M(x_1, \ldots, x_n)$ n'est pas
très clair ! Nous utiliserons plutôt l'idée de la démonstration du principe de
réflexion pour établir directement les résultats d'existence dont nous avons
besoin.

\section{Construction de catégories de schémas}
\label{section-categories-schemes}

\noindent
Nous allons expliquer comment appliquer ceci pour produire, à partir d'un
ensemble initial de schémas, une « petite » catégorie de schémas stable par une
liste d'opérations naturelles. Avant cela, introduisons la taille d'un schéma.
Étant donné un schéma $S$, nous définissons
$$
\text{size}(S) = \max(\aleph_0, \kappa_1, \kappa_2),
$$
où nous définissons les nombres cardinaux $\kappa_1$ et $\kappa_2$ comme suit :
\begin{enumerate}
\item Nous prenons pour $\kappa_1$ la cardinalité de l'ensemble des ouverts
affines de $S$.
\item Pour $\kappa_2$, nous prenons la borne supérieure des cardinalités des
$\Gamma(U, \mathcal{O}_S)$, lorsque $U \subset S$ est un ouvert affine.
\end{enumerate}

\begin{lemma}
\label{lemma-bounded-size}
Pour tout cardinal $\kappa$, il existe un ensemble $A$ tel que chaque élément
de $A$ soit un schéma et que, pour tout schéma $S$ vérifiant
$\text{size}(S) \leq \kappa$, il existe un élément $X \in A$ tel que
$X \cong S$ (isomorphisme de schémas).
\end{lemma}

\begin{proof}
La démonstration est omise. Indication : réfléchir à la manière dont tout schéma est isomorphe à un
schéma obtenu par recollement de schémas affines.
\end{proof}

\noindent
Nous notons $Bound$ la fonction qui associe à chaque cardinal $\kappa$
\begin{equation}
\label{equation-bound}
Bound(\kappa) = \max\{\kappa^{\aleph_0}, \kappa^+\}.
\end{equation}
Nous pourrions faire croître cette fonction beaucoup plus rapidement ; par
exemple, nous pourrions poser $Bound(\kappa) = \kappa^\kappa$, et le résultat
ci-dessous resterait vrai. Pour tout ordinal $\alpha$, nous notons
$\Sch_\alpha$ la sous-catégorie pleine de la catégorie des schémas dont les
objets sont des éléments de $V_\alpha$. Voici le résultat que nous allons
démontrer.

\begin{lemma}
\label{lemma-construct-category}
Avec les notations $\text{size}$, $Bound$ et $\Sch_\alpha$ ci-dessus, soit
$S_0$ un ensemble de schémas. Il existe un ordinal limite $\alpha$ possédant
les propriétés suivantes :
\begin{enumerate}
\item
\label{item-inclusion}
Nous avons $S_0 \subset V_\alpha$ ; autrement dit,
$S_0 \subset \Ob(\Sch_\alpha)$.
\item
\label{item-bounded}
Pour tout $S \in \Ob(\Sch_\alpha)$ et tout schéma $T$ tel que
$\text{size}(T) \leq Bound(\text{size}(S))$, il existe un schéma
$S' \in \Ob(\Sch_\alpha)$ tel que $T \cong S'$.
\item
\label{item-limit}
Pour toute catégorie d'indices dénombrable\footnote{L'ensemble des objets
et les ensembles de morphismes sont tous dénombrables. On peut en fait
démontrer le lemme en remplaçant $\aleph_0$ par un cardinal quelconque dans
(3) et (4).} $\mathcal{I}$ et tout foncteur
$F : \mathcal{I} \to \Sch_\alpha$, la limite $\lim_\mathcal{I} F$ existe dans
$\Sch_\alpha$ si et seulement si elle existe dans $\Sch$ et, de plus, dans ce
cas, le morphisme naturel entre elles est un isomorphisme.
\item
\label{item-colimit}
Pour toute catégorie d'indices dénombrable $\mathcal{I}$ et tout foncteur
$F : \mathcal{I} \to \Sch_\alpha$, la colimite $\colim_\mathcal{I} F$ existe
dans $\Sch_\alpha$ si et seulement si elle existe dans $\Sch$ et, de plus,
dans ce cas, le morphisme naturel entre elles est un isomorphisme.
\end{enumerate}
\end{lemma}

\begin{proof}
Nous définissons, par récurrence transfinie, une fonction $f$ qui associe un
ordinal à chaque ordinal comme suit. Posons $f(0) = 0$. Étant donné
$f(\alpha)$, nous définissons $f(\alpha + 1)$ comme le plus petit ordinal
$\beta$ pour lequel les conditions suivantes sont satisfaites :
\begin{enumerate}
\item Nous avons $\alpha + 1 \leq \beta$ et $f(\alpha) \leq \beta$.
\item Soient $S \in \Ob(\Sch_{f(\alpha)})$ et $T$ un schéma tel que
$\text{size}(T) \leq Bound(\text{size}(S))$. Il existe alors un schéma
$S' \in \Ob(\Sch_\beta)$ tel que $T \cong S'$.
\item Pour toute catégorie d'indices dénombrable $\mathcal{I}$ et tout
foncteur $F : \mathcal{I} \to \Sch_{f(\alpha)}$, si la limite
$\lim_\mathcal{I} F$ ou la colimite $\colim_\mathcal{I} F$ existe dans
$\Sch$, alors elle est isomorphe à un schéma de $\Sch_\beta$.
\end{enumerate}
Pour voir que $\beta$ existe, raisonnons comme suit. Puisque
$\Ob(\Sch_{f(\alpha)})$ est un ensemble, nous voyons que
$\kappa =
\sup_{S \in \Ob(\Sch_{f(\alpha)})} Bound(\text{size}(S))$
existe et est un cardinal. Soit $A$ un ensemble de schémas obtenu à partir de
$\kappa$ comme dans le lemme \ref{lemma-bounded-size}. Il existe un ensemble
$CountCat$ de catégories dénombrables tel que toute catégorie dénombrable soit
isomorphe à un élément de $CountCat$. Nous pouvons donc supposer, au point (3)
ci-dessus, que $\mathcal{I}$ est un élément de $CountCat$. Cela signifie que
les paires $(\mathcal{I}, F)$ du point (3) parcourent un ensemble. Il existe
donc un ensemble $B$ dont les éléments sont des schémas tel que, pour toute
paire $(\mathcal{I}, F)$ comme au point (3), si la limite ou la colimite existe,
elle soit isomorphe à un élément de $B$. Par conséquent, si nous choisissons un
$\beta$ quelconque tel que $A \cup B \subset V_\beta$ et
$\beta > \max\{\alpha + 1, f(\alpha)\}$, alors (1)--(3) sont satisfaites.
Comme toute collection non vide d'ordinaux possède un plus petit élément, nous
voyons que $f(\alpha + 1)$ est bien définie. Enfin, si $\alpha$ est un ordinal
limite, nous posons
$f(\alpha) = \sup_{\alpha' < \alpha} f(\alpha')$.

\medskip\noindent
Choisissons $\beta_0$ tel que $S_0 \subset V_{\beta_0}$. Par construction,
$f(\beta) \geq \beta$, et nous voyons que
$S_0 \subset V_{f(\beta_0)}$ également. De plus, comme $f$ est croissante au
sens large, nous voyons que $S_0 \subset V_{f(\beta)}$ est vrai pour tout
$\beta \geq \beta_0$. Choisissons ensuite un ordinal quelconque
$\beta_1 > \beta_0$ de cofinalité
$\text{cf}(\beta_1) > \omega = \aleph_0$. C'est possible, car la cofinalité des
ordinaux devient arbitrairement grande ; voir la proposition
\ref{proposition-exist-ordinals-large-cofinality}. Nous affirmons que
$\alpha = f(\beta_1)$ est une solution au problème posé dans le lemme.

\medskip\noindent
La première propriété du lemme découle de notre choix de
$\beta_1 > \beta_0$ ci-dessus.

\medskip\noindent
Puisque $\beta_1$ est un ordinal limite (sa cofinalité étant infinie), nous
obtenons $f(\beta_1) = \sup_{\beta < \beta_1} f(\beta)$. Ainsi,
$\{f(\beta) \mid \beta < \beta_1\} \subset f(\beta_1)$ est un sous-ensemble
cofinal. Nous voyons donc que
$$
V_\alpha = V_{f(\beta_1)} = \bigcup\nolimits_{\beta < \beta_1} V_{f(\beta)}.
$$
Soit maintenant $S \in \Ob(\Sch_\alpha)$. Nous définissons $\beta(S)$ comme le
plus petit ordinal $\beta$ tel que
$S \in \Ob(\Sch_{f(\beta)})$. D'après ce qui précède, nous voyons que
$\beta(S) < \beta_1$ toujours. Puisque
$\Ob(\Sch_{f(\beta + 1)}) \subset
\Ob(\Sch_\alpha)$, la construction de $f$ ci-dessus montre que la deuxième
propriété du lemme est satisfaite.

\medskip\noindent
Supposons que $\{S_1, S_2, \ldots\} \subset \Ob(\Sch_\alpha)$ soit une
collection dénombrable. Considérons la fonction
$\omega \to \beta_1$, $n \mapsto \beta(S_n)$. Comme la cofinalité de $\beta_1$
est $> \omega$, l'image de cette fonction ne peut pas être un sous-ensemble
cofinal. Il existe donc un $\beta < \beta_1$ tel que
$\{S_1, S_2, \ldots\} \subset \Ob(\Sch_{f(\beta)})$. Il s'ensuit que tout
foncteur $F : \mathcal{I} \to \Sch_\alpha$ se factorise par l'une des
sous-catégories $\Sch_{f(\beta)}$. Ainsi, s'il existe un schéma $X$ qui est la
colimite ou la limite du diagramme $F$, la construction de $f$ montre que $X$
est isomorphe à un objet de $\Sch_{f(\beta + 1)}$, qui est une sous-catégorie
de $\Sch_\alpha$. Cela démontre les deux dernières assertions du lemme.
\end{proof}

\begin{remark}
\label{remark-how-to-use-reflection}
Le lemme ci-dessus peut également être démontré à l'aide du principe de
réflexion. Il faut toutefois être prudent. Supposons en effet que la phrase
$\phi_{scheme}(X)$ exprime la propriété « $X$ est un schéma » ; que signifie
alors la formule $\phi_{scheme}^{V_\alpha}(X)$ ? Il est vrai que le principe de
réflexion affirme que nous pouvons trouver $\alpha$ tel que, pour tout
$X \in V_\alpha$, nous ayons
$\phi_{scheme}(X) \leftrightarrow \phi_{scheme}^{V_\alpha}(X)$,
mais cela ne sert absolument à rien. C'est seulement en combinant deux énoncés
de ce type que quelque chose d'intéressant se produit. Supposons par exemple
que $\phi_{red}(X, Y)$ exprime la propriété « $X$, $Y$ sont des schémas et $Y$
est la réduction de $X$ » (voir Schémas, définition
\ref{schemes-definition-reduced-induced-scheme}). Supposons que nous appliquions
le principe de réflexion à la paire de formules
$\phi_1(X, Y) = \phi_{red}(X, Y)$,
$\phi_2(X) = \exists Y, \phi_1(X, Y)$. Il est alors facile de voir que tout
$\alpha$ produit par le principe de réflexion possède la propriété suivante :
pour tout $X \in \Ob(\Sch_\alpha)$, la réduction de $X$ est également un objet
de $\Sch_\alpha$ (vérification laissée en exercice).
\end{remark}

\begin{lemma}
\label{lemma-bound-affine}
Soit $S$ un schéma affine. Soit $R = \Gamma(S, \mathcal{O}_S)$. Alors la taille
de $S$ est égale à $\max\{ \aleph_0, |R|\}$.
\end{lemma}

\begin{proof}
Il y a au plus $\max\{|R|, \aleph_0\}$ ouverts affines de $\Spec(R)$. Cela est
clair, car tout ouvert affine $U \subset \Spec(R)$ est une union finie
d'ouverts principaux $D(f_1) \cup \ldots \cup D(f_n)$ ; le nombre d'ouverts
affines est donc au plus
$\sup_n |R|^n = \max\{|R|, \aleph_0\}$, voir
\cite[chap. I, 10.13]{Kunen}. D'autre part, nous voyons que
$\Gamma(U, \mathcal{O}) \subset R_{f_1} \times \ldots \times R_{f_n}$ et donc
$|\Gamma(U, \mathcal{O})| \leq
\max\{\aleph_0, |R_{f_1}|, \ldots, |R_{f_n}|\}$. Il suffit donc de démontrer
que $|R_f| \leq \max\{\aleph_0, |R|\}$, ce que nous omettons.
\end{proof}

\begin{lemma}
\label{lemma-bound-size}
Soit $S$ un schéma. Soit $S = \bigcup_{i \in I} S_i$ un recouvrement ouvert.
Alors
$\text{size}(S) \leq \max\{|I|, \sup_i\{\text{size}(S_i)\}\}$.
\end{lemma}

\begin{proof}
Soit $U \subset S$ un ouvert affine quelconque. Comme $U$ est quasi-compact,
il existe un nombre fini d'éléments $i_1, \ldots, i_n \in I$ et des ouverts
affines $U_i \subset U \cap S_i$ tels que
$U = U_1 \cup U_2 \cup \ldots \cup U_n$. Ainsi
$$
|\Gamma(U, \mathcal{O}_U)|
\leq
|\Gamma(U_1, \mathcal{O})|
\otimes
\ldots
\otimes
|\Gamma(U_n, \mathcal{O})|
\leq \sup\nolimits_i\{\text{size}(S_i)\}
$$
De plus, cela montre que la cardinalité de l'ensemble des ouverts affines de
$S$ est inférieure ou égale à celle de l'ensemble
$$
\coprod_{n \in \omega}
\coprod_{i_1, \ldots, i_n \in I}
\{\text{ouverts affines de }S_{i_1}\}
\times
\ldots
\times
\{\text{ouverts affines de }S_{i_n}\}.
$$
La cardinalité de chacun des ensembles figurant dans l'union disjointe est
majorée par $\sup_i\{\text{size}(S_i)\}$. Celle de l'ensemble d'indices est
au plus égale à $\max\{|I|, \aleph_0\}$, voir
\cite[chap. I, 10.13]{Kunen}. Par conséquent, d'après
\cite[lemme 5.8]{Jech}, la cardinalité du coproduit est au plus
$$
\max\{\aleph_0, |I|\}
\otimes \sup_i\{\text{size}(S_i)\}
$$
Le lemme en résulte.
\end{proof}

\begin{lemma}
\label{lemma-bound-size-fibre-product}
Soient $f : X \to S$, $g : Y \to S$ des morphismes de schémas. Alors nous
avons
$\text{size}(X \times_S Y) \leq \max\{\text{size}(X), \text{size}(Y)\}$.
\end{lemma}

\begin{proof}
Soit $S = \bigcup_{k \in K} S_k$ un recouvrement ouvert affine. Soient
$X = \bigcup_{i \in I} U_i$, $Y = \bigcup_{j \in J} V_j$ des recouvrements
ouverts affines, où $I$, $J$ ont des cardinalités
$\leq \text{size}(X), \text{size}(Y)$. Pour chaque $i \in I$, il existe un
ensemble fini $K_i$ d'éléments $k \in K$ tel que
$f(U_i) \subset \bigcup_{k \in K_i} S_k$. Pour chaque $j \in J$, il existe un
ensemble fini $K_j$ d'éléments $k \in K$ tel que
$g(V_j) \subset \bigcup_{k \in K_j} S_k$. Par conséquent, $f(X), g(Y)$ sont
contenus dans $S' = \bigcup_{k \in K'} S_k$, où
$K' = \bigcup_{i \in I} K_i \cup \bigcup_{j \in J} K_j$. Remarquons que la
cardinalité de $K'$ est au plus
$\max\{\aleph_0, |I|, |J|\}$. En appliquant le lemme
\ref{lemma-bound-size}, nous voyons qu'il suffit de démontrer que
$\text{size}(f^{-1}(S_k) \times_{S_k} g^{-1}(S_k))
\leq \max\{\text{size}(X), \text{size}(Y)\}$ pour $k \in K'$. Autrement dit,
nous pouvons supposer que $S$ est affine.

\medskip\noindent
Supposons $S$ affine. Soient
$X = \bigcup_{i \in I} U_i$, $Y = \bigcup_{j \in J} V_j$ des recouvrements
ouverts affines, où $I$, $J$ ont des cardinalités
$\leq \text{size}(X), \text{size}(Y)$. De nouveau, d'après le lemme
\ref{lemma-bound-size}, il suffit de démontrer le lemme pour les produits
$U_i \times_S V_j$. D'après le lemme \ref{lemma-bound-affine}, nous voyons
qu'il suffit de montrer que
$$
|A \otimes_C B| \leq \max\{\aleph_0, |A|, |B|\}.
$$
Nous omettons la démonstration de cette inégalité.
\end{proof}

\begin{lemma}
\label{lemma-bound-finite-type}
Soit $S$ un schéma. Soit $f : X \to S$ un morphisme localement de type fini,
où $X$ est quasi-compact. Alors
$\text{size}(X) \leq \text{size}(S)$.
\end{lemma}

\begin{proof}
Nous pouvons trouver un recouvrement ouvert affine fini
$X = \bigcup_{i = 1, \ldots n} U_i$ tel que chaque $U_i$ s'envoie dans un
ouvert affine $S_i$ de $S$. Ainsi, d'après le lemme
\ref{lemma-bound-size}, nous nous ramenons au cas où $S$ et $X$ sont tous deux
affines. Dans ce cas, le lemme \ref{lemma-bound-affine} montre qu'il suffit de
prouver
$$
|A[x_1, \ldots, x_n]| \leq \max\{\aleph_0, |A|\}.
$$
Nous omettons la démonstration de cette inégalité.
\end{proof}

\noindent
Dans Algèbre, lemme \ref{algebra-lemma-epimorphism-cardinality}, nous
montrerons que si $A \to B$ est un épimorphisme d'anneaux, alors
$|B| \leq \max(|A|, \aleph_0)$. L'analogue pour les schémas est le lemme
suivant.

\begin{lemma}
\label{lemma-bound-monomorphism}
Soit $f : X \to Y$ un monomorphisme de schémas. Si au moins l'une des
propriétés suivantes est satisfaite, alors
$\text{size}(X) \leq \text{size}(Y)$ :
\begin{enumerate}
\item $f$ est quasi-compact,
\item $f$ est localement de présentation finie,
\item ajouter ici d'autres cas selon les besoins.
\end{enumerate}
Mais cette borne n'est pas valable pour les monomorphismes localement de type
fini.
\end{lemma}

\begin{proof}
Soit $Y = \bigcup_{j \in J} V_j$ un recouvrement ouvert affine de $Y$ tel que
$|J| \leq \text{size}(Y)$. D'après le lemme \ref{lemma-bound-size}, il suffit
de borner la taille de l'image réciproque de $V_j$ dans $X$. Nous nous ramenons
donc au cas où $Y$ est affine, disons $Y = \Spec(B)$. Pour tout ouvert affine
$\Spec(A) \subset X$, nous avons
$|A| \leq \max(|B|, \aleph_0) = \text{size}(Y)$ ; voir la remarque ci-dessus
et le lemme \ref{lemma-bound-affine}. Il suffit donc de montrer que $X$ possède
au plus $\text{size}(Y)$ ouverts affines. Cela est clair si $X$ est
quasi-compact, d'où le cas (1). Dans le cas (2), le nombre de classes
d'isomorphisme de $B$-algèbres $A$ susceptibles d'apparaître est borné par
$\text{size}(B)$, car chaque $A$ est de type fini sur $B$, donc isomorphe à une
algèbre $B[x_1, \ldots, x_n]/(f_1, \ldots, f_m)$ pour certains $n, m$ et
$f_j \in B[x_1, \ldots, x_n]$. Cependant, comme $X \to Y$ est un
monomorphisme, il existe au plus un morphisme $\Spec(A) \to X$ au-dessus de
$Y = \Spec(B)$. Le nombre d'ouverts affines de $X$ est donc borné par le nombre
de ces classes d'isomorphisme.

\medskip\noindent
Pour démontrer la dernière assertion du lemme, considérons l'anneau
$B = \prod_{n \in \mathbf{N}} \mathbf{F}_2$ et posons $Y = \Spec(B)$. Pour
chaque ultrafiltre $\mathcal{U}$ sur $\mathbf{N}$, nous obtenons un idéal
maximal $\mathfrak m_\mathcal{U}$ de corps résiduel $\mathbf{F}_2$ ;
l'application $B \to \mathbf{F}_2$ envoie l'élément $(x_n)$ sur
$\lim_\mathcal{U} x_n$. Les détails sont omis. Le morphisme de schémas
$X = \coprod_\mathcal{U} \Spec(\mathbf{F}_2) \to Y$ est un monomorphisme,
puisque tous les points sont distincts. Cependant, la cardinalité de l'ensemble
des sous-schémas ouverts affines de $X$ est égale à celle de l'ensemble des
ultrafiltres sur $\mathbf{N}$, qui vaut $2^{2^{\aleph_0}}$. Nous concluons grâce
à $|B| = 2^{\aleph_0} < 2^{2^{\aleph_0}}$.
\end{proof}
\begin{lemma}
\label{lemma-what-is-in-it}
Soit $\alpha$ un ordinal comme dans le lemme \ref{lemma-construct-category}
ci-dessus. La catégorie $\Sch_\alpha$ vérifie les propriétés suivantes :
\begin{enumerate}
\item Si $X, Y, S \in \Ob(\Sch_\alpha)$, alors, pour tous morphismes
$f : X \to S$, $g : Y \to S$, le produit fibré $X \times_S Y$ dans
$\Sch_\alpha$ existe et est un produit fibré dans la catégorie des schémas.
\item Étant donnée toute famille au plus dénombrable $S_1, S_2, \ldots$
d'éléments de $\Ob(\Sch_\alpha)$, le coproduit $\coprod_i S_i$ existe dans
$\Ob(\Sch_\alpha)$ et est un coproduit dans la catégorie des schémas.
\item Pour tout $S \in \Ob(\Sch_\alpha)$ et toute immersion ouverte
$U \to S$, il existe un $V \in \Ob(\Sch_\alpha)$ tel que $V \cong U$.
\item Pour tout $S \in \Ob(\Sch_\alpha)$ et toute immersion fermée
$T \to S$, il existe un $S' \in \Ob(\Sch_\alpha)$ tel que $S' \cong T$.
\item Pour tout $S \in \Ob(\Sch_\alpha)$ et tout morphisme de type fini
$T \to S$, il existe un $S' \in \Ob(\Sch_\alpha)$ tel que $S' \cong T$.
\item Supposons que $S$ soit un schéma possédant un recouvrement ouvert
$S = \bigcup_{i \in I} S_i$ tel qu'il existe un $T \in \Ob(\Sch_\alpha)$
avec (a) $\text{size}(S_i) \leq \text{size}(T)^{\aleph_0}$ pour tout
$i \in I$, et (b) $|I| \leq \text{size}(T)^{\aleph_0}$.
Alors $S$ est isomorphe à un objet de $\Sch_\alpha$.
\item Pour tout $S \in \Ob(\Sch_\alpha)$ et tout morphisme $f : T \to S$
localement de type fini tel que $T$ puisse être recouvert par au plus
$\text{size}(S)^{\aleph_0}$ ouverts affines, il existe un
$S' \in \Ob(\Sch_\alpha)$ tel que $S' \cong T$.
Par exemple, cela vaut si $T$ peut être recouvert par au plus
$|\mathbf{R}| = 2^{\aleph_0} = \aleph_0^{\aleph_0}$ ouverts affines.
\item Pour tout $S \in \Ob(\Sch_\alpha)$ et tout monomorphisme $T \to S$
qui est soit localement de présentation finie, soit quasi-compact, il existe
un $S' \in \Ob(\Sch_\alpha)$ tel que $S' \cong T$.
\item Supposons que $T \in \Ob(\Sch_\alpha)$ soit affine. Écrivons
$R = \Gamma(T, \mathcal{O}_T)$. Alors chacun des schémas suivants est
isomorphe à un schéma de $\Sch_\alpha$ :
\begin{enumerate}
\item Pour tout idéal $I \subset R$, si $R^* = \lim_n R/I^n$ est le
complété, le schéma $\Spec(R^*)$.
\item Pour toute $R$-algèbre de type fini $R'$, le schéma $\Spec(R')$.
\item Pour toute localisation $S^{-1}R$, le schéma $\Spec(S^{-1}R)$.
\item Pour tout idéal premier $\mathfrak p \subset R$, le schéma
$\Spec(\overline{\kappa(\mathfrak p)})$.
\item Pour tout sous-anneau $R' \subset R$, le schéma $\Spec(R')$.
\item Tout schéma de type fini sur un anneau de cardinalité au plus
$|R|^{\aleph_0}$.
\item Et ainsi de suite.
\end{enumerate}
\end{enumerate}
\end{lemma}

\begin{proof}
Les assertions (1) et (2) découlent directement des définitions.
L'assertion (3) découle du fait que la taille d'un sous-schéma ouvert $U$ de
$S$ est clairement inférieure ou égale à celle de $S$.
L'assertion (4) découle de (5).
L'assertion (5) découle de (7).
L'assertion (6) découle du fait que la taille de $S$ est
$\leq \max\{|I|, \sup_i \text{size}(S_i)\} \leq \text{size}(T)^{\aleph_0}$
d'après le lemme \ref{lemma-bound-size}. L'assertion (7) découle de (6).
En effet, pour tout ouvert affine $V \subset T$, nous avons
$\text{size}(V) \leq \text{size}(S)$ d'après le
lemme \ref{lemma-bound-finite-type}.
Nous voyons donc que (6) s'applique dans la situation de (7).
La partie (8) découle du lemme \ref{lemma-bound-monomorphism}.

\medskip\noindent
L'assertion (9) se traduit, via le lemme \ref{lemma-bound-affine}, par une
borne supérieure pour la cardinalité des anneaux
$R^*$, $S^{-1}R$, $\overline{\kappa(\mathfrak p)}$, $R'$, etc.
Le cas le plus intéressant est peut-être l'anneau $R^*$. En tant qu'ensemble,
il est l'image d'une application surjective $R^{\mathbf{N}} \to R^*$.
Puisque $|R^{\mathbf{N}}| = |R|^{\aleph_0}$, la borne voulue résulte de ce que
nous avons choisi $Bound(\kappa)$ supérieur ou égal à $\kappa^{\aleph_0}$.
Ouf ! (La cardinalité de la clôture algébrique d'un corps
est la même que celle du corps, ou bien elle vaut $\aleph_0$.)
\end{proof}

\begin{remark}
\label{remark-what-is-not-in-it}
Soit $R$ un anneau. Supposons que nous considérions l'anneau
$\prod_{\mathfrak p \in \Spec(R)} \kappa(\mathfrak p)$.
La cardinalité de cet anneau est bornée par $|R|^{2^{|R|}}$, mais n'est pas
bornée par $|R|^{\aleph_0}$ en général.
Par exemple, si $R = \mathbf{C}[x]$, elle n'est pas bornée par
$|R|^{\aleph_0}$, et si $R = \prod_{n \in \mathbf{N}} \mathbf{F}_2$,
elle n'est pas bornée par $|R|^{|R|}$.
Ainsi, le « Et ainsi de suite » du lemme \ref{lemma-what-is-in-it} ci-dessus
doit être pris avec prudence. Bien entendu, s'il devient un jour nécessaire
de considérer ces anneaux dans des arguments relatifs à la cohomologie
fppf/\'etale, nous pourrons remplacer la fonction $Bound$ ci-dessus par la
fonction $\kappa \mapsto \kappa^{2^\kappa}$.
\end{remark}

\noindent
Dans le lemme suivant, nous utilisons la notion de recouvrement fpqc introduite
dans Topologies, section \ref{topologies-section-fpqc}.

\begin{lemma}
\label{lemma-bound-by-covering}
Soit $f : X \to Y$ un morphisme de schémas. Supposons qu'il existe un
recouvrement fpqc $\{g_j : Y_j \to Y\}_{j \in J}$ tel que $g_j$ se factorise
par $f$. Alors $\text{size}(Y) \leq \text{size}(X)$.
\end{lemma}

\begin{proof}
Soit $V \subset Y$ un ouvert affine. Par définition, il existe un entier
$n \geq 0$, une application $a : \{1, \ldots, n\} \to J$ et des ouverts
affines $V_i \subset Y_{a(i)}$ tels que
$V = g_{a(1)}(V_1) \cup \ldots \cup g_{a(n)}(V_n)$.
Notons $h_j : Y_j \to X$ un morphisme tel que $f \circ h_j = g_j$.
Alors $h_{a(1)}(V_1) \cup \ldots \cup h_{a(n)}(V_n)$ est une partie
quasi-compacte de $f^{-1}(V)$. Nous pouvons donc trouver un ouvert
quasi-compact $W \subset f^{-1}(V)$ qui contient $h_{a(i)}(V_i)$ pour
$i = 1, \ldots, n$. En particulier, $V = f(W)$.

\medskip\noindent
D'une part, cela montre que la cardinalité de l'ensemble des ouverts affines
de $Y$ est au plus celle de l'ensemble $S$ des ouverts quasi-compacts de $X$.
Puisque tout ouvert quasi-compact de $X$ est une réunion finie d'ouverts
affines, nous voyons que la cardinalité de cet ensemble est au plus
$\sup |S|^n = \max(\aleph_0, |S|)$.
D'autre part, nous avons
$\mathcal{O}_Y(V) \subset \prod_{i = 1, \ldots, n} \mathcal{O}_{Y_{a(i)}}(V_i)$
car $\{V_i \to V\}$ est un recouvrement fpqc. Par conséquent,
$\mathcal{O}_Y(V) \subset \mathcal{O}_X(W)$ puisque $V_i \to V$ se factorise
par $W$. Comme $W$ possède lui aussi un recouvrement fini par des ouverts
affines de $X$, nous concluons que $|\mathcal{O}_Y(V)|$ est borné par la
taille de $X$. Le lemme découle maintenant de la définition de la taille d'un
schéma.
\end{proof}

\noindent
Dans le lemme suivant, nous utilisons la notion de recouvrement fppf introduite
dans Topologies, section \ref{topologies-section-fppf}.

\begin{lemma}
\label{lemma-bound-fppf-covering}
Soit $\{f_i : X_i \to X\}_{i \in I}$ un recouvrement fppf d'un schéma.
Il existe un recouvrement fppf $\{W_j \to X\}_{j \in J}$ qui raffine
$\{X_i \to X\}_{i \in I}$ et tel que
$\text{size}(\coprod W_j) \leq \text{size}(X)$.
\end{lemma}

\begin{proof}
Choisissons un recouvrement ouvert affine $X = \bigcup_{a \in A} U_a$ avec
$|A| \leq \text{size}(X)$. Pour chaque $a$, nous pouvons choisir une partie
finie $I_a \subset I$ et, pour $i \in I_a$, un ouvert quasi-compact
$W_{a, i} \subset X_i$ tels que
$U_a = \bigcup_{i \in I_a} f_i(W_{a, i})$.
Alors $\text{size}(W_{a, i}) \leq \text{size}(X)$ d'après le
lemme \ref{lemma-bound-finite-type}.
Nous concluons que
$$
\text{size}(\coprod_a \coprod_{i \in I_a} W_{i, a}) \leq \text{size}(X)
$$
d'après le lemme \ref{lemma-bound-size}.
\end{proof}





\section{Ensembles munis d'une action de groupe}
\label{section-sets-with-group-action}

\noindent
Soit $G$ un groupe. Nous notons $G\textit{-Sets}$ la « grande » catégorie
des $G$-ensembles. Pour tout ordinal $\alpha$, nous notons
$G\textit{-Sets}_\alpha$ la sous-catégorie pleine de $G\textit{-Sets}$ dont
les objets appartiennent à $V_\alpha$. Comme notion de taille d'un
$G$-ensemble, nous prenons
$\text{size}(S) = \max\{\aleph_0, |G|, |S|\}$ (où $|G|$ et $|S|$ sont les
cardinalités des ensembles sous-jacents). Comme ci-dessus, nous utilisons la
fonction $Bound(\kappa) = \kappa^{\aleph_0}$.

\begin{lemma}
\label{lemma-sets-with-group-action}
Avec les notations $G$, $G\textit{-Sets}_\alpha$, $\text{size}$ et $Bound$
ci-dessus, soit $S_0$ un ensemble de $G$-ensembles.
Il existe un ordinal limite $\alpha$ ayant les propriétés suivantes :
\begin{enumerate}
\item Nous avons $S_0 \cup \{{}_GG\} \subset \Ob(G\textit{-Sets}_\alpha)$.
\item Soit $S \in \Ob(G\textit{-Sets}_\alpha)$.\par\noindent
Pour tout $G$-ensemble $T$ tel
que $\text{size}(T) \leq Bound(\text{size}(S))$, il existe un
$S' \in \Ob(G\textit{-Sets}_\alpha)$ qui est isomorphe à $T$.
\item Soit $\mathcal{I}$ une catégorie d'indices dénombrable. Pour tout foncteur
$F : \mathcal{I} \to G\textit{-Sets}_\alpha$, la limite
$\lim_\mathcal{I} F$ et la colimite $\colim_\mathcal{I} F$ existent dans
$G\textit{-Sets}_\alpha$ et coïncident avec ceux calculés dans $G\textit{-Sets}$.
\end{enumerate}
\end{lemma}

\begin{proof}
Omis. La démonstration est analogue à celle du
lemme \ref{lemma-construct-category} ci-dessus, mais plus facile.
\end{proof}

\begin{lemma}
\label{lemma-what-is-in-it-G-sets}
Soit $\alpha$ un ordinal comme dans le
lemme \ref{lemma-sets-with-group-action} ci-dessus. La catégorie
$G\textit{-Sets}_\alpha$ vérifie les propriétés suivantes :
\begin{enumerate}
\item Le $G$-ensemble ${}_GG$ est un objet de $G\textit{-Sets}_\alpha$.
\item Les (co)produits, les produits fibrés et les sommes amalgamées existent
dans $G\textit{-Sets}_\alpha$ et coïncident avec ceux calculés dans
$G\textit{-Sets}$.
\item Étant donné un objet $U$ de $G\textit{-Sets}_\alpha$, toute partie
$G$-stable $O \subset U$ est isomorphe à un objet de
$G\textit{-Sets}_\alpha$.
\end{enumerate}
\end{lemma}

\begin{proof}
Omis.
\end{proof}

\section{Recouvrements d'un site}
\label{section-coverings-site}

\noindent
Supposons que $\mathcal{C}$ soit une catégorie (comme dans Catégories,
définition \ref{categories-definition-category}) et que
$\text{Cov}(\mathcal{C})$ soit une classe propre de recouvrements satisfaisant
les propriétés (1), (2) et (3) de Sites,
définition \ref{sites-definition-site}. Nous les rappelons ici :
\begin{enumerate}
\item Si $V \to U$ est un isomorphisme, alors $\{V \to U\} \in
\text{Cov}(\mathcal{C})$.
\item Si $\{U_i \to U\}_{i\in I} \in \text{Cov}(\mathcal{C})$ et si, pour
chaque $i$, nous avons
$\{V_{ij} \to U_i\}_{j\in J_i} \in \text{Cov}(\mathcal{C})$, alors
$\{V_{ij} \to U\}_{i \in I, j\in J_i} \in \text{Cov}(\mathcal{C})$.
\item Si $\{U_i \to U\}_{i\in I}\in \text{Cov}(\mathcal{C})$ et si
$V \to U$ est un morphisme de $\mathcal{C}$, alors $U_i \times_U V$ existe
pour tout $i$ et
$\{U_i \times_U V \to V \}_{i\in I} \in \text{Cov}(\mathcal{C})$.
\end{enumerate}
Pour un ordinal $\alpha$, nous posons
$\text{Cov}(\mathcal{C})_\alpha = \text{Cov}(\mathcal{C}) \cap V_\alpha$.
Étant donnés un ordinal $\alpha$ et un cardinal $\kappa$, nous définissons
$\text{Cov}(\mathcal{C})_{\kappa, \alpha}$ comme l'ensemble des éléments
$\mathcal{U} =
\{\varphi_i : U_i \to U\}_{i\in I} \in \text{Cov}(\mathcal{C})_\alpha$
tels que $|I| \leq \kappa$.

\medskip\noindent
Rappelons la notion suivante ; voir Sites,
définition \ref{sites-definition-combinatorial-tautological}.
Deux familles de morphismes, $\{\varphi_i : U_i \to U\}_{i\in I}$ et
$\{\psi_j : W_j \to U\}_{j\in J}$, ayant la même cible dans $\mathcal{C}$,
sont dites {\it combinatoirement équivalentes} s'il existe des applications
$\alpha : I \to J$ et $\beta : J\to I$ telles que
$\varphi_i = \psi_{\alpha(i)}$ et $\psi_j = \varphi_{\beta(j)}$.
Cela définit une relation d'équivalence sur les familles de morphismes ayant
une cible fixée.

\begin{lemma}
\label{lemma-coverings-site}
Avec les notations ci-dessus, soit
$\text{Cov}_0 \subset \text{Cov}(\mathcal{C})$ un ensemble contenu dans
$\text{Cov}(\mathcal{C})$.
Il existe un cardinal $\kappa$ et un ordinal limite $\alpha$ ayant les
propriétés suivantes :
\begin{enumerate}
\item Nous avons
$\text{Cov}_0 \subset \text{Cov}(\mathcal{C})_{\kappa, \alpha}$.
\item L'ensemble de recouvrements
$\text{Cov}(\mathcal{C})_{\kappa, \alpha}$ satisfait les propriétés (1), (2)
et (3) de Sites, définition \ref{sites-definition-site} (voir ci-dessus).
Autrement dit, $(\mathcal{C}, \text{Cov}(\mathcal{C})_{\kappa, \alpha})$
est un site.
\item Tout recouvrement de $\text{Cov}(\mathcal{C})$ est combinatoirement
équivalent à un recouvrement de
$\text{Cov}(\mathcal{C})_{\kappa, \alpha}$.
\end{enumerate}
\end{lemma}

\begin{proof}
Pour le démontrer, considérons d'abord l'ensemble $\mathcal{S}$ de tous les
ensembles de morphismes de $\mathcal{C}$ ayant une cible fixée.
Autrement dit, un élément de $\mathcal{S}$ est une partie $T$ de
$\text{Arrows}(\mathcal{C})$ telle que tous les éléments de $T$ aient la même
cible. Étant donnée une famille
$\mathcal{U} = \{\varphi_i : U_i \to U\}_{i\in I}$ de morphismes ayant une
cible fixée, nous définissons
$Supp(\mathcal{U}) = \{ \varphi \in \text{Arrows}(\mathcal{C})
\mid \exists i\in I, \varphi = \varphi_i\}$.
Remarquons que deux familles
$\mathcal{U} =  \{\varphi_i : U_i \to U\}_{i\in I}$ et
$\mathcal{V} = \{V_j \to V\}_{j \in J}$ sont combinatoirement équivalentes
si et seulement si $Supp(\mathcal{U}) = Supp(\mathcal{V})$.
Ensuite, nous définissons $\mathcal{S}_\tau \subset \mathcal{S}$ comme la
partie
$\mathcal{S}_\tau = \{ T \in \mathcal{S} \mid
\exists\ \mathcal{U} \in \text{Cov}(\mathcal{C}) \ T = Supp(\mathcal{U})\}$.
Pour tout élément $T \in \mathcal{S}_\tau$, posons $\beta(T)$ égal au plus
petit ordinal $\beta$ tel qu'il existe un
$\mathcal{U} \in \text{Cov}(\mathcal{C})_\beta$ tel que
$T = \text{Supp}(\mathcal{U})$. Enfin, posons
$\beta_0 = \sup_{T \in S_\tau} \beta(T)$.
Il en résulte que tout $\mathcal{U} \in \text{Cov}(\mathcal{C})$ est
combinatoirement équivalent à un élément de
$\text{Cov}(\mathcal{C})_{\beta_0}$.

\medskip\noindent
Soit $\kappa$ le maximum de $\aleph_0$, de la cardinalité
$|\text{Arrows}(\mathcal{C})|$,
$$
\sup\nolimits_{\{U_i \to U\}_{i\in I} \in \text{Cov}(\mathcal{C})_{\beta_0}}
|I|,
\quad\text{et}\quad
\sup\nolimits_{\{U_i \to U\}_{i\in I} \in \text{Cov}_0} |I|.
$$
Puisque $\kappa$ est un cardinal infini, nous avons
$\kappa \otimes \kappa = \kappa$. Remarquons qu'évidemment
$\text{Cov}(\mathcal{C})_{\beta_0} =
\text{Cov}(\mathcal{C})_{\kappa, \beta_0}$.

\medskip\noindent
Nous définissons, par récurrence transfinie, une fonction $f$ qui associe un
ordinal à tout ordinal de la manière suivante. Posons $f(0) = 0$.
Étant donné $f(\alpha)$, nous définissons $f(\alpha + 1)$ comme le plus petit
ordinal $\beta$ tel que les conditions suivantes soient satisfaites :
\begin{enumerate}
\item Nous avons $\alpha + 1 \leq \beta$ et $f(\alpha) \leq \beta$.
\item Si $\{U_i \to U\}_{i\in I}
\in \text{Cov}(\mathcal{C})_{\kappa, f(\alpha)}$ et si, pour tout $i$, nous
avons
$\{W_{ij} \to U_i\}_{j\in J_i}
\in \text{Cov}(\mathcal{C})_{\kappa, f(\alpha)}$,
alors
$\{W_{ij} \to U\}_{i \in I, j\in J_i}
\in \text{Cov}(\mathcal{C})_{\kappa, \beta}$.
\item Si $\{U_i \to U\}_{i\in I}
\in \text{Cov}(\mathcal{C})_{\kappa, \alpha}$ et si $W \to U$ est un
morphisme de $\mathcal{C}$, alors
$\{U_i \times_U W \to W \}_{i\in I}
\in \text{Cov}(\mathcal{C})_{\kappa, \beta}$.
\end{enumerate}
Pour voir que $\beta$ existe, remarquons que la collection de tous les
recouvrements $\{W_{ij} \to U\}$ et $\{U_i \times_U W \to W \}$ qui
interviennent dans (2) et (3) forme clairement un ensemble. Il existe donc un
ordinal $\beta$ tel que $V_\beta$ contienne tous ces recouvrements. En outre,
l'ensemble d'indices du recouvrement $\{W_{ij} \to U\}$ a pour cardinalité
$\sum_{i \in I} |J_i| \leq \kappa \otimes \kappa = \kappa$, et ces
recouvrements sont donc contenus dans
$\text{Cov}(\mathcal{C})_{\kappa, \beta}$.
Puisque toute collection non vide d'ordinaux possède un plus petit élément,
nous voyons que $f(\alpha + 1)$ est bien défini. Enfin, si $\alpha$ est un
ordinal limite, nous posons
$f(\alpha) = \sup_{\alpha' < \alpha} f(\alpha')$.

\medskip\noindent
Choisissons un ordinal $\beta_1$ tel que
$\text{Arrows}(\mathcal{C}) \subset V_{\beta_1}$,
$\text{Cov}_0 \subset V_{\beta_0}$ et $\beta_1 \geq \beta_0$.
Par construction, $f(\beta_1) \geq \beta_1$, et nous voyons que les mêmes
propriétés sont vérifiées pour $V_{f(\beta_1)}$. De plus, comme $f$ est
croissante au sens large, cela reste vrai pour tout $\beta \geq \beta_1$.
Choisissons ensuite un ordinal $\beta_2 > \beta_1$ de cofinalité
$\text{cf}(\beta_2) > \kappa$. C'est possible puisque la cofinalité des
ordinaux devient arbitrairement grande ; voir la
proposition \ref{proposition-exist-ordinals-large-cofinality}.
Nous affirmons que la paire formée par $\kappa$ et $\alpha = f(\beta_2)$ est
une solution au problème posé dans le lemme.

\medskip\noindent
La première et la troisième propriétés du lemme découlent de nos choix de
$\kappa$, $\beta_2 > \beta_1 > \beta_0$ ci-dessus.

\medskip\noindent
Puisque $\beta_2$ est un ordinal limite (car sa cofinalité est infinie), nous
obtenons $f(\beta_2) = \sup_{\beta < \beta_2} f(\beta)$.
Ainsi, $\{f(\beta) \mid \beta < \beta_2\} \subset f(\beta_2)$ est une partie
cofinale. Nous voyons donc que
$$
V_\alpha = V_{f(\beta_2)} = \bigcup\nolimits_{\beta < \beta_2} V_{f(\beta)}.
$$
Maintenant, soit $\mathcal{U} \in \text{Cov}_{\kappa, \alpha}$.
Nous définissons $\beta(\mathcal{U})$ comme le plus petit ordinal $\beta$ tel
que $\mathcal{U} \in \text{Cov}_{\kappa, f(\beta)}$. D'après ce qui précède,
nous voyons que $\beta(\mathcal{U}) < \beta_2$ dans tous les cas.

\medskip\noindent
Nous devons montrer que les propriétés (1), (2) et (3) qui définissent un site
sont vérifiées pour la paire
$(\mathcal{C}, \text{Cov}_{\kappa, \alpha})$.
La première l'est parce que, par notre choix de $\beta_2$, toutes les flèches
de $\mathcal{C}$ appartiennent à $V_{f(\beta_2)}$.
Pour la troisième, étant donné un recouvrement
$\mathcal{U} = \{U_i \to U\}_{i \in I}
\in \text{Cov}(\mathcal{C})_{\kappa, \alpha}$, nous utilisons le fait que
$\beta(\mathcal{U}) < \beta_2$ ; tout changement de base de $\mathcal{U}$
appartient alors, par construction de $f$, à
$\text{Cov}(\mathcal{C})_{\kappa, f(\beta + 1)}$, et donc à
$\text{Cov}(\mathcal{C})_{\kappa, \alpha}$.

\medskip\noindent
Enfin, pour la deuxième condition, supposons que
$\{U_i \to U\}_{i\in I}
\in \text{Cov}(\mathcal{C})_{\kappa, f(\alpha)}$ et que, pour chaque $i$,
nous ayons
$\mathcal{W}_i = \{W_{ij} \to U_i\}_{j\in J_i}
\in \text{Cov}(\mathcal{C})_{\kappa, f(\alpha)}$.
Considérons la fonction
$I \to \beta_2$, $i \mapsto \beta(\mathcal{W}_i)$. Puisque la cofinalité
de $\beta_2$ est $> \kappa \geq |I|$, l'image de cette fonction ne peut être
une partie cofinale. Il existe donc un $\beta < \beta_1$ tel que
$\mathcal{W}_i \in \text{Cov}_{\kappa, f(\beta)}$ pour tout $i \in I$.
Il en résulte que le recouvrement
$\{W_{ij} \to U\}_{i\in I, j \in J_i}$ est un élément de
$\text{Cov}(\mathcal{C})_{\kappa, f(\beta + 1)}
\subset \text{Cov}(\mathcal{C})_{\kappa, \alpha}$, comme souhaité.
\end{proof}

\begin{remark}
\label{remark-better}
Il est probable que, pour un certain ordinal limite $\alpha$, l'ensemble de
recouvrements $\text{Cov}(\mathcal{C})_\alpha$ satisfasse les conditions du
lemme. C'est après tout ce qu'une application du principe de réflexion semble
donner (aux réserves près décrites à la fin de la
section \ref{section-reflection-principle} et dans la
remarque \ref{remark-how-to-use-reflection}).
\end{remark}








\section{Catégories abéliennes et injectifs}
\label{section-abelian-categories-injectives}

\noindent
Le lemme suivant s'applique à la catégorie des modules sur un faisceau
d'anneaux sur un site.

\begin{lemma}
\label{lemma-abelian-injectives}
Supposons donnée une grande catégorie $\mathcal{A}$ (voir Catégories,
remarque \ref{categories-remark-big-categories}).
Supposons que $\mathcal{A}$ soit abélienne et possède suffisamment d'injectifs.
Voir Homologie, définitions \ref{homology-definition-abelian-category}
et \ref{homology-definition-enough-injectives}.
Alors, pour tout ensemble donné d'objets $\{A_s\}_{s\in S}$ de
$\mathcal{A}$, il existe une sous-catégorie abélienne
$\mathcal{A}' \subset \mathcal{A}$ ayant les propriétés suivantes :
\begin{enumerate}
\item le foncteur d'inclusion $\mathcal{A}' \to \mathcal{A}$ est exact,
\item $\Ob(\mathcal{A}')$ est un ensemble,
\item $\Ob(\mathcal{A}')$ contient $A_s$ pour tout $s \in S$,
\item $\mathcal{A}'$ possède suffisamment d'injectifs, et
\item un objet de $\mathcal{A}'$ est injectif si et seulement s'il est un
objet injectif de $\mathcal{A}$.
\end{enumerate}
\end{lemma}

\begin{proof}
Omis.
\end{proof}









\begin{multicols}{2}[\section{Autres chapitres}]
\noindent
Préliminaires
\begin{enumerate}
\item \hyperref[introduction-section-phantom]{Introduction}
\item \hyperref[conventions-section-phantom]{Conventions}
\item \hyperref[sets-section-phantom]{Théorie des ensembles}
\item \hyperref[categories-section-phantom]{Catégories}
\item \hyperref[topology-section-phantom]{Topologie}
\item \hyperref[sheaves-section-phantom]{Faisceaux sur les espaces}
\item \hyperref[sites-section-phantom]{Sites et faisceaux}
\item \hyperref[stacks-section-phantom]{Champs}
\item \hyperref[fields-section-phantom]{Corps}
\item \hyperref[algebra-section-phantom]{Algèbre commutative}
\item \hyperref[brauer-section-phantom]{Groupes de Brauer}
\item \hyperref[homology-section-phantom]{Algèbre homologique}
\item \hyperref[derived-section-phantom]{Catégories dérivées}
\item \hyperref[simplicial-section-phantom]{Méthodes simpliciales}
\item \hyperref[more-algebra-section-phantom]{Compléments d'algèbre}
\item \hyperref[smoothing-section-phantom]{Lissage des morphismes d'anneaux}
\item \hyperref[modules-section-phantom]{Faisceaux de modules}
\item \hyperref[sites-modules-section-phantom]{Modules sur les sites}
\item \hyperref[injectives-section-phantom]{Objets injectifs}
\item \hyperref[cohomology-section-phantom]{Cohomologie des faisceaux}
\item \hyperref[sites-cohomology-section-phantom]{Cohomologie sur les sites}
\item \hyperref[dga-section-phantom]{Algèbre différentielle graduée}
\item \hyperref[dpa-section-phantom]{Algèbre à puissances divisées}
\item \hyperref[sdga-section-phantom]{Faisceaux différentiels gradués}
\item \hyperref[hypercovering-section-phantom]{Hyperrecouvrements}
\end{enumerate}
Schémas
\begin{enumerate}
\setcounter{enumi}{25}
\item \hyperref[schemes-section-phantom]{Schémas}
\item \hyperref[constructions-section-phantom]{Constructions de schémas}
\item \hyperref[properties-section-phantom]{Propriétés des schémas}
\item \hyperref[morphisms-section-phantom]{Morphismes de schémas}
\item \hyperref[coherent-section-phantom]{Cohomologie des schémas}
\item \hyperref[divisors-section-phantom]{Diviseurs}
\item \hyperref[limits-section-phantom]{Limites de schémas}
\item \hyperref[varieties-section-phantom]{Variétés}
\item \hyperref[topologies-section-phantom]{Topologies sur les schémas}
\item \hyperref[descent-section-phantom]{Descente}
\item \hyperref[perfect-section-phantom]{Catégories dérivées de schémas}
\item \hyperref[more-morphisms-section-phantom]{Compléments sur les morphismes}
\item \hyperref[flat-section-phantom]{Compléments sur la platitude}
\item \hyperref[groupoids-section-phantom]{Schémas en groupoïdes}
\item \hyperref[more-groupoids-section-phantom]{Compléments sur les schémas en groupoïdes}
\item \hyperref[etale-section-phantom]{Morphismes étales de schémas}
\end{enumerate}
Thèmes de théorie des schémas
\begin{enumerate}
\setcounter{enumi}{41}
\item \hyperref[chow-section-phantom]{Homologie de Chow}
\item \hyperref[intersection-section-phantom]{Théorie de l'intersection}
\item \hyperref[pic-section-phantom]{Schémas de Picard des courbes}
\item \hyperref[weil-section-phantom]{Théories cohomologiques de Weil}
\item \hyperref[adequate-section-phantom]{Modules adéquats}
\item \hyperref[dualizing-section-phantom]{Complexes dualisants}
\item \hyperref[duality-section-phantom]{Dualité pour les schémas}
\item \hyperref[discriminant-section-phantom]{Discriminants et différentes}
\item \hyperref[derham-section-phantom]{Cohomologie de de Rham}
\item \hyperref[local-cohomology-section-phantom]{Cohomologie locale}
\item \hyperref[algebraization-section-phantom]{Géométrie algébrique et formelle}
\item \hyperref[curves-section-phantom]{Courbes algébriques}
\item \hyperref[resolve-section-phantom]{Résolution des surfaces}
\item \hyperref[models-section-phantom]{Réduction semi-stable}
\item \hyperref[functors-section-phantom]{Foncteurs et morphismes}
\item \hyperref[equiv-section-phantom]{Catégories dérivées de variétés}
\item \hyperref[pione-section-phantom]{Groupes fondamentaux de schémas}
\item \hyperref[etale-cohomology-section-phantom]{Cohomologie étale}
\item \hyperref[crystalline-section-phantom]{Cohomologie cristalline}
\item \hyperref[proetale-section-phantom]{Cohomologie pro-étale}
\item \hyperref[relative-cycles-section-phantom]{Cycles relatifs}
\item \hyperref[more-etale-section-phantom]{Compléments de cohomologie étale}
\item \hyperref[trace-section-phantom]{La formule des traces}
\end{enumerate}
Espaces algébriques
\begin{enumerate}
\setcounter{enumi}{64}
\item \hyperref[spaces-section-phantom]{Espaces algébriques}
\item \hyperref[spaces-properties-section-phantom]{Propriétés des espaces algébriques}
\item \hyperref[spaces-morphisms-section-phantom]{Morphismes d'espaces algébriques}
\item \hyperref[decent-spaces-section-phantom]{Espaces algébriques décents}
\item \hyperref[spaces-cohomology-section-phantom]{Cohomologie des espaces algébriques}
\item \hyperref[spaces-limits-section-phantom]{Limites d'espaces algébriques}
\item \hyperref[spaces-divisors-section-phantom]{Diviseurs sur les espaces algébriques}
\item \hyperref[spaces-over-fields-section-phantom]{Espaces algébriques sur des corps}
\item \hyperref[spaces-topologies-section-phantom]{Topologies sur les espaces algébriques}
\item \hyperref[spaces-descent-section-phantom]{Descente et espaces algébriques}
\item \hyperref[spaces-perfect-section-phantom]{Catégories dérivées d'espaces}
\item \hyperref[spaces-more-morphisms-section-phantom]{Compléments sur les morphismes d'espaces}
\item \hyperref[spaces-flat-section-phantom]{Platitude sur les espaces algébriques}
\item \hyperref[spaces-groupoids-section-phantom]{Groupoïdes dans les espaces algébriques}
\item \hyperref[spaces-more-groupoids-section-phantom]{Compléments sur les groupoïdes dans les espaces}
\item \hyperref[bootstrap-section-phantom]{Amorçage}
\item \hyperref[spaces-pushouts-section-phantom]{Sommes amalgamées d'espaces algébriques}
\end{enumerate}
Thèmes de géométrie
\begin{enumerate}
\setcounter{enumi}{81}
\item \hyperref[spaces-chow-section-phantom]{Groupes de Chow des espaces}
\item \hyperref[groupoids-quotients-section-phantom]{Quotients de groupoïdes}
\item \hyperref[spaces-more-cohomology-section-phantom]{Compléments sur la cohomologie des espaces}
\item \hyperref[spaces-simplicial-section-phantom]{Espaces simpliciaux}
\item \hyperref[spaces-duality-section-phantom]{Dualité pour les espaces}
\item \hyperref[formal-spaces-section-phantom]{Espaces algébriques formels}
\item \hyperref[restricted-section-phantom]{Algébrisation des espaces formels}
\item \hyperref[spaces-resolve-section-phantom]{Résolution des surfaces revisitée}
\end{enumerate}
Théorie des déformations
\begin{enumerate}
\setcounter{enumi}{89}
\item \hyperref[formal-defos-section-phantom]{Théorie formelle des déformations}
\item \hyperref[defos-section-phantom]{Théorie des déformations}
\item \hyperref[cotangent-section-phantom]{Le complexe cotangent}
\item \hyperref[examples-defos-section-phantom]{Problèmes de déformation}
\end{enumerate}
Champs algébriques
\begin{enumerate}
\setcounter{enumi}{93}
\item \hyperref[algebraic-section-phantom]{Champs algébriques}
\item \hyperref[examples-stacks-section-phantom]{Exemples de champs}
\item \hyperref[stacks-sheaves-section-phantom]{Faisceaux sur les champs algébriques}
\item \hyperref[criteria-section-phantom]{Critères de représentabilité}
\item \hyperref[artin-section-phantom]{Axiomes d'Artin}
\item \hyperref[quot-section-phantom]{Espaces de Quot et de Hilbert}
\item \hyperref[stacks-properties-section-phantom]{Propriétés des champs algébriques}
\item \hyperref[stacks-morphisms-section-phantom]{Morphismes de champs algébriques}
\item \hyperref[stacks-limits-section-phantom]{Limites de champs algébriques}
\item \hyperref[stacks-cohomology-section-phantom]{Cohomologie des champs algébriques}
\item \hyperref[stacks-perfect-section-phantom]{Catégories dérivées de champs}
\item \hyperref[stacks-introduction-section-phantom]{Introduction aux champs algébriques}
\item \hyperref[stacks-more-morphisms-section-phantom]{Compléments sur les morphismes de champs}
\item \hyperref[stacks-geometry-section-phantom]{La géométrie des champs}
\end{enumerate}
Thèmes de théorie des espaces de modules
\begin{enumerate}
\setcounter{enumi}{107}
\item \hyperref[moduli-section-phantom]{Champs de modules}
\item \hyperref[moduli-curves-section-phantom]{Espaces de modules de courbes}
\end{enumerate}
Divers
\begin{enumerate}
\setcounter{enumi}{109}
\item \hyperref[examples-section-phantom]{Exemples}
\item \hyperref[exercises-section-phantom]{Exercices}
\item \hyperref[guide-section-phantom]{Guide bibliographique}
\item \hyperref[desirables-section-phantom]{Desiderata}
\item \hyperref[coding-section-phantom]{Style de programmation}
\item \hyperref[obsolete-section-phantom]{Obsolète}
\item \hyperref[fdl-section-phantom]{Licence de documentation libre GNU}
\item \hyperref[index-section-phantom]{Index généré automatiquement}
\end{enumerate}
\end{multicols}


\bibliography{my}
\bibliographystyle{amsalpha}

\end{document}
